\chapter{Introducción}
\label{chap:introduccion}
En el presente informe se desarrolla el Trabajo Práctico Nº 1, cuyo objetivo principal es estudiar y aplicar métodos
de resolución de sistemas lineales e invariantes en el tiempo (SLIT). Para ello, se analizan tres circuitos eléctricos
que representan este tipo de sistemas, abordando su comportamiento tanto desde el punto de vista teórico como práctico.

A lo largo del trabajo se aplican los métodos vistos en clase para resolver los sistemas propuestos, obteniendo sus
respuestas en frecuencia y analizando las características principales de cada circuito. Además, se realizan
simulaciones con el fin de verificar los resultados teóricos y comparar el comportamiento esperado con el obtenido
mediante herramientas computacionales.

Finalmente, se implementan los circuitos de manera experimental y se miden sus respuestas en frecuencia, lo que permite
contrastar los resultados prácticos con los valores calculados y simulados. De esta forma, el trabajo busca integrar
el análisis matemático, la simulación y la implementación física de sistemas SLIT, fortaleciendo la comprensión de su
funcionamiento y sus aplicaciones en circuitos eléctricos.
